\sectionkicker{Source-backed technical notes}
\subsection*{能力拡張と同時に増えた制御面: Technical Notes}
本欄は記事本文で圧縮した一次資料上の情報を、比較・再検証しやすい形で整理したものである。時系列、確認済みの主張、留意点、一次資料URLを掲載する。
\smallskip
\noindent{\small\textit{Technical Notesでは、一次資料から確認した公開・提供・研究上の事実を記載する。数値・能力評価や提供元・著者による比較表現は、独立再現済みの結果ではなく帰属claimとして扱う。}}\par
\medskip
\subsection*{Theme at a glance}
\begin{center}
\footnotesize
\begin{tabularx}{\linewidth}{@{}p{0.20\linewidth}p{0.10\linewidth}p{0.14\linewidth}X@{}}
\toprule
資料 & 位置づけ & 種別 & 時系列 \\
\midrule
Many-shot jailbreaking & 主要資料 & 安全性研究 & 2024-04-02             \\
Anthropic interpretability research (Mapping the Mind / Golden Gate Claude) & 主要資料 & 公式情報 & 2024-05-21, 2024-05-23             \\
OpenAI Model Spec & 主要資料 & 公式情報 & 2024-05-08             \\
Claude 3 → Claude 3.5 Sonnet & 補足資料 & モデル更新 & 2024-03-04, 2024-06-21             \\
Gemini 1.5 / Gemini 1.5 Flash & 補足資料 & モデル更新 & 2024-02-15, 2024-05-14             \\
\bottomrule
\end{tabularx}
\end{center}
\normalsize

\Needspace{0.38\textheight}
\begin{technicalnote}{Many-shot jailbreaking}{主要資料}
\begingroup
% reader-facing Technical Notes break policy
\widowpenalty=10000
\clubpenalty=10000
\displaywidowpenalty=10000
\begin{tabularx}{\linewidth}{@{}>{\bfseries}p{0.22\linewidth}X@{}}
組織 & Anthropic \\
種別 & 安全性研究 \\
時系列 & 2024-04-02 \\
\end{tabularx}
\smallskip
\begin{minipage}{\linewidth}
% reader-facing Technical Notes event-only fact/source tail group
{\bfseries 一次資料から整理したtechnical points}
\begin{itemize}[leftmargin=1.5em,itemsep=0.35em]
\item \textbf{一次情報で確認できる事実}: Anthropicの一次資料により、Many-shot jailbreakingについて2024-04-02にResearch公開を確認した。 Many-shot jailbreakingは、長いcontextへ多数のfaux dialogue demonstrationsを並べ、最後の危険なrequestに対するsafety behaviorを崩すlong-context attackとしてAnthropicが報告した。shot数が増えるほどharmful responseが生じやすくなる傾向と、既存jailbreak技法との組合せで必要prompt lengthを短縮できる結果を示しているが、攻撃成功率等は著者報告として扱う。
\end{itemize}
\begin{minipage}{\linewidth}
% reader-facing Technical Notes generic boundary/source tail group
\begin{samepage}
{\bfseries 一次資料}
\begingroup\sloppy
\Urlmuskip=0mu plus 2mu\relax
\begin{itemize}[leftmargin=1.5em,itemsep=0.25em]
\item {\scriptsize\url{https://www.anthropic.com/research/many-shot-jailbreaking}}
\end{itemize}
\endgroup
\end{samepage}
\end{minipage}
\end{minipage}
\endgroup
\end{technicalnote}
\medskip

\Needspace{0.38\textheight}
\begin{technicalnote}{Anthropic interpretability research (Mapping the Mind / Golden Gate Claude)}{主要資料}
\begingroup
% reader-facing Technical Notes break policy
\widowpenalty=10000
\clubpenalty=10000
\displaywidowpenalty=10000
\begin{tabularx}{\linewidth}{@{}>{\bfseries}p{0.22\linewidth}X@{}}
組織 & Anthropic \\
種別 & 公式情報 \\
時系列 & 2024-05-21, 2024-05-23 \\
\end{tabularx}
\smallskip
\begin{minipage}{\linewidth}
% reader-facing Technical Notes event-only fact/source tail group
{\bfseries 一次資料から整理したtechnical points}
\begin{itemize}[leftmargin=1.5em,itemsep=0.35em]
\item \textbf{一次情報で確認できる事実}: Anthropicの一次資料により、Anthropic interpretability research (Mapping the Mind / Golden Gate Claude)について2024-05-21にResearch公開、2024-05-23にResearch公開を確認した。 Anthropicはdictionary learningをClaude 3 Sonnetのmiddle layerへscaleし、neuron activation patternsをhuman-interpretable featuresへ分解してmillions of featuresを抽出したと報告した。featuresはmultilingual/multimodalなconceptにも反応し、feature activationを人工的に増減するとmodel behaviorが変化することをGolden Gate Claude demoで示した。一方、抽出featureは全conceptの一部で、featureが属するcircuitsの理解は未解決と明記している。
\end{itemize}
\begin{minipage}{\linewidth}
% reader-facing Technical Notes generic boundary/source tail group
\begin{samepage}
{\bfseries 一次資料}
\begingroup\sloppy
\Urlmuskip=0mu plus 2mu\relax
\begin{itemize}[leftmargin=1.5em,itemsep=0.25em]
\item {\scriptsize\url{https://www.anthropic.com/news/golden-gate-claude}}
\item {\scriptsize\url{https://www.anthropic.com/research/mapping-mind-language-model}}
\end{itemize}
\endgroup
\end{samepage}
\end{minipage}
\end{minipage}
\endgroup
\end{technicalnote}
\medskip

\Needspace{0.38\textheight}
\begin{technicalnote}{OpenAI Model Spec}{主要資料}
\begingroup
% reader-facing Technical Notes break policy
\widowpenalty=10000
\clubpenalty=10000
\displaywidowpenalty=10000
\begin{tabularx}{\linewidth}{@{}>{\bfseries}p{0.22\linewidth}X@{}}
組織 & OpenAI \\
種別 & 公式情報 \\
時系列 & 2024-05-08 \\
\end{tabularx}
\smallskip
\begin{minipage}{\linewidth}
% reader-facing Technical Notes event-only fact/source tail group
{\bfseries 一次資料から整理したtechnical points}
\begin{itemize}[leftmargin=1.5em,itemsep=0.35em]
\item \textbf{一次情報で確認できる事実}: OpenAIの一次資料により、OpenAI Model Specについて2024-05-08にSpec公開を確認した。 2024年5月8日のModel Spec初版は、望ましいmodel behaviorをObjectives・Rules・Default behaviorsの3層で整理し、Rulesで安全性・合法性に関わる境界を定めつつ、Defaultsはその境界内でdeveloper/userが上書き可能なstarting pointとして扱う構造を示した。OpenAIは当時、この文書をRLHFに関わるresearchers/data labelers向けのguidelineとして用いる意図を示す一方、現行形そのものはまだ学習に直接使っていないと明記しており、仕様公開と実モデル挙動を同一視しない。
\end{itemize}
\begin{minipage}{\linewidth}
% reader-facing Technical Notes generic boundary/source tail group
\begin{samepage}
{\bfseries 一次資料}
\begingroup\sloppy
\Urlmuskip=0mu plus 2mu\relax
\begin{itemize}[leftmargin=1.5em,itemsep=0.25em]
\item {\scriptsize\url{https://openai.com/index/introducing-the-model-spec}}
\end{itemize}
\endgroup
\end{samepage}
\end{minipage}
\end{minipage}
\endgroup
\end{technicalnote}
\medskip

\Needspace{0.38\textheight}
\begin{technicalnote}{Claude 3 → Claude 3.5 Sonnet}{補足資料}
\begingroup
% reader-facing Technical Notes break policy
\widowpenalty=10000
\clubpenalty=10000
\displaywidowpenalty=10000
\begin{tabularx}{\linewidth}{@{}>{\bfseries}p{0.22\linewidth}X@{}}
組織 & Anthropic \\
種別 & モデル更新 \\
時系列 & 2024-03-04, 2024-06-21 \\
\end{tabularx}
\smallskip
\begin{minipage}{\linewidth}
% reader-facing Technical Notes event-only fact/source tail group
{\bfseries 一次資料から整理したtechnical points}
\begin{itemize}[leftmargin=1.5em,itemsep=0.35em]
\item \textbf{一次情報で確認できる事実}: Anthropicの一次資料により、Claude 3 → Claude 3.5 Sonnetについて2024-03-04にモデル公開、2024-06-21にモデル公開を確認した。 対象event近傍の一次資料から 200K context を確認できる。
\end{itemize}
\begin{minipage}{\linewidth}
% reader-facing Technical Notes generic boundary/source tail group
\begin{samepage}
{\bfseries 一次資料}
\begingroup\sloppy
\Urlmuskip=0mu plus 2mu\relax
\begin{itemize}[leftmargin=1.5em,itemsep=0.25em]
\item {\scriptsize\url{https://www.anthropic.com/news/claude-3-5-sonnet}}
\item {\scriptsize\url{https://www.anthropic.com/news/claude-3-family}}
\end{itemize}
\endgroup
\end{samepage}
\end{minipage}
\end{minipage}
\endgroup
\end{technicalnote}
\medskip

\Needspace{0.38\textheight}
\begin{technicalnote}{Gemini 1.5 / Gemini 1.5 Flash}{補足資料}
\begingroup
% reader-facing Technical Notes break policy
\widowpenalty=10000
\clubpenalty=10000
\displaywidowpenalty=10000
\begin{tabularx}{\linewidth}{@{}>{\bfseries}p{0.22\linewidth}X@{}}
組織 & Google \\
種別 & モデル更新 \\
時系列 & 2024-02-15, 2024-05-14 \\
\end{tabularx}
\smallskip
\begin{minipage}{\linewidth}
% reader-facing Technical Notes event-only fact/source tail group
{\bfseries 一次資料から整理したtechnical points}
\begin{itemize}[leftmargin=1.5em,itemsep=0.35em]
\item \textbf{一次情報で確認できる事実}: Googleの一次資料により、Gemini 1.5 / Gemini 1.5 Flashについて2024-02-15にモデルPreview、2024-05-14にモデル発表を確認した。 対象event近傍の一次資料から Mixture-of-Experts (MoE) を確認できる。
\end{itemize}
\begin{minipage}{\linewidth}
% reader-facing Technical Notes generic boundary/source tail group
\begin{samepage}
{\bfseries 一次資料}
\begingroup\sloppy
\Urlmuskip=0mu plus 2mu\relax
\begin{itemize}[leftmargin=1.5em,itemsep=0.25em]
\item {\scriptsize\url{https://blog.google/innovation-and-ai/products/google-gemini-next-generation-model-february-2024/}}
\item {\scriptsize\url{https://blog.google/innovation-and-ai/technology/developers-tools/gemini-gemma-developer-updates-may-2024/}}
\end{itemize}
\endgroup
\end{samepage}
\end{minipage}
\end{minipage}
\endgroup
\end{technicalnote}
\medskip
